\chapter*{АННОТАЦИЯ}
\addcontentsline{toc}{chapter}{АННОТАЦИЯ}

Работа посвящена методам автоматизированной генерации системно-динамических моделей из текстовых описаний с использованием больших языковых моделей. Цель исследования --- разработать и реализовать подход, при котором вероятностная генерация дополняется типизированным промежуточным представлением, детерминированными построителями и проверками корректности.

Предложен конвейер, объединяющий краткое описание задачи, план типизированных компонентов, каноническое представление уравнений, проверку единиц измерения и топологии, ограниченное исправление и контрольную симуляцию. Для явно заданных дифференциальных уравнений применяется построитель, сохраняющий коэффициенты и факторизованную форму исходных выражений. Генерация Python-кода непосредственно языковой моделью не используется.

Реализованы автономный модуль генерации, REST API, backend и веб-интерфейс. Поддерживаются контекстные беседы, версии моделей, симуляция, сценарии, анализ чувствительности, калибровка и ограниченный импорт Vensim. На наборе SDXorg импортированы и симулированы 63 модели; 62 из 148 моделей с эталонными траекториями совпали с ними при неизменных допусках. Это соответствует 98,4\% внутри поддерживаемого подмножества и 41,9\% всего корпуса.

Проведены оценка локальной Qwen3-8B-Q4\_K\_M на 70 acceptance-генерациях и контролируемое сравнение трёх вариантов генерации на 216 случаях замороженного набора. Полный конвейер обеспечил исполняемость 94,44\% моделей; семантические и поведенческие проверки прошли 87,50\% и 84,72\% соответственно. Ограниченное исправление увеличило долю поведенческих успехов на 16,67 п.п., однако порог 85\% формально не достигнут. Сопоставимое сравнение с 30B остаётся отдельным экспериментом.

Научная новизна заключается в применении типизированных промежуточных представлений и детерминированных ограничений для контроля LLM-генерации системно-динамических моделей. Практическая значимость состоит в создании воспроизводимой платформы для исследования таких методов и быстрого прототипирования моделей. Контролируемое сравнение выполнено для Ministral-8B и локальной Qwen3-8B; независимая экспертная оценка остаётся следующим необходимым этапом.

\textbf{Ключевые слова:} системная динамика, большие языковые модели, типизированное промежуточное представление, детерминированная генерация, проверка единиц измерения, топология моделей, Vensim, PySD, локальные языковые модели.
